% chapter-3.tex
% Chapter 3 — Methodology
%
% This chapter describes the planned experimental methodology. Citations
% use the citation keys of the project Zotero library (exported to
% references.bib). Items that depend on decisions not yet finalized
% (hardware details, software versions, model selection, resolutions,
% durations, repetitions) are marked with visible \todo{...} placeholders
% to be completed once determined; % TODO comments mark citation gaps. No
% experimental results are reported in this chapter.

\providecommand{\todo}[1]{\textbf{[TODO: #1]}}

\chapter{Methodology}
\label{chap:methodology}

\section{Experimental Platform}
\label{sec:platform}

The primary experimental platform is an NVIDIA Jetson Orin Nano, chosen
as a representative embedded module for edge computer vision
\cite{scalcon2024, alqahtani2026}. The following characteristics of the
device are recorded and held constant throughout the experiments:

\begin{itemize}
  \item Jetson module/model: \todo{SPECIFY: exact module and developer kit
        version}
  \item RAM capacity: \todo{SPECIFY}
  \item Storage: \todo{SPECIFY}
  \item Power mode: \todo{SPECIFY: fixed power mode for all experiments.
        Published Orin-family evaluations distinguish power modes when
        reporting frame rates \cite{scalcon2024}, so this setting is part
        of the documented experimental condition.}
  \item Cooling configuration: \todo{SPECIFY. Sustained inference load
        raises device temperature \cite{zequn2026}, so cooling is fixed
        and documented.}
  \item Connected camera / input device: \todo{SPECIFY}
  \item Operating temperature, where relevant: \todo{SPECIFY: how thermal
        conditions are monitored}
\end{itemize}

All experiments are executed on this single device, with nodes of the
ROS~2 configurations running as separate operating-system processes on
the same host. Running all processes on one device isolates the effect of
the ROS~2 architecture itself from network effects, which are outside the
scope of this thesis.

\section{Software Environment}
\label{sec:software-environment}

The complete software environment is documented and held constant:

\begin{itemize}
  \item Operating system: \todo{SPECIFY: Linux distribution and version}
  \item JetPack version: \todo{SPECIFY}
  \item ROS~2 distribution: \todo{SPECIFY: one distribution, fixed for all
        experiments}
  \item DDS implementation: \todo{SPECIFY: one implementation, fixed for
        all experiments, selected through the rmw layer}
  \item QoS configuration: \todo{SPECIFY: one fixed configuration, stated
        explicitly, including reliability, durability, and history
        policies. QoS settings measurably affect latency
        \cite{fernandez2020, kronauer2021} and must therefore be part of
        the documented experimental condition.}
  \item Python version: \todo{SPECIFY}
  \item Package versions (inference runtime, image libraries, monitoring
        tools): \todo{SPECIFY}
\end{itemize}

A non-containerized deployment is used for the primary experiments. The
choices above are treated as controlled experimental conditions rather
than variables under investigation; each is documented so that the
experiments can be reproduced (Section~\ref{sec:reproducibility}).

\section{Computer-Vision Pipeline}
\label{sec:pipeline}

The workload is an object-detection pipeline consisting of the following
logical stages:

\begin{enumerate}
  \item \textbf{Image acquisition:} obtaining each input frame.
  \item \textbf{Preprocessing:} converting the frame to the format expected
        by the model (e.g., resizing, color conversion, normalization, and
        tensor layout). This stage is individually significant: for object
        detection, preprocessing contributes a substantial portion of
        total runtime, especially at higher resolutions
        \cite{szegedi2026}.
  \item \textbf{Inference:} executing the neural-network model.
  \item \textbf{Postprocessing:} converting raw model output into
        application-level results (e.g., bounding boxes, confidence
        filtering, non-maximum suppression). Like preprocessing, this
        stage is non-negligible for object detection \cite{szegedi2026}.
  \item \textbf{Output:} producing the final result of the pipeline.
\end{enumerate}

The neural-network model is selected according to the criteria of
Section~\ref{sec:workload}, drawing on published evidence that real-time
object detection is feasible on the Orin Nano with standard one-stage
detector families \cite{scalcon2024, alqahtani2026, szegedi2026,zequn2026}:
\todo{SPECIFY: final model, weights version, precision, and inference
runtime once selected.} The same model and inference configuration are
used in every experimental configuration.

\section{Monolithic Implementation}
\label{sec:monolithic}

The baseline implementation executes the complete pipeline within a single
Python process:

\begin{quote}
\texttt{Input $\rightarrow$ Preprocessing $\rightarrow$ Inference
$\rightarrow$ Postprocessing $\rightarrow$ Output}
\end{quote}

All stages execute sequentially within one process. Data passes between
stages as in-memory data structures, with no serialization or
inter-process transport. This implementation serves as the performance
baseline against which the ROS~2 implementations are compared. It
corresponds to the architectural extreme in which the application takes
full advantage of the fact that all pipeline data can remain in a single
address space.

\section{ROS 2 Implementation}
\label{sec:ros2-implementation}

The same logical stages are implemented as ROS~2 nodes using rclpy, each
running as a separate operating-system process on the same host, connected
in a chain by topics carrying image messages:

\begin{quote}
\texttt{Node A $\rightarrow$ Node B $\rightarrow$ Node C $\rightarrow$
Node D}
\end{quote}

Each node publishes and subscribes with the fixed QoS configuration, and
executor settings are likewise fixed and identical across nodes, since
the executor is inseparable from node performance
\cite{macenski2023, casini2026}. The mapping of processing stages to
nodes preserves functional equivalence with the monolithic implementation:
given the same input stream and model, both implementations produce
equivalent outputs.
\todo{SPECIFY: the exact stage-to-node mapping for each configuration once
finalized, the executor type used, and how functional equivalence is
verified (e.g., comparing outputs on a fixed input sequence).}

This implementation corresponds to the ordinary, recommended-by-default
way of structuring a ROS~2 application as separate nodes
\cite{macenski2022}, before any of the mitigations described in the
related literature---composition, intra-process communication, or
shared-memory transport \cite{macenski2023, kamiya2025}---are applied.
Studying this baseline in isolation is deliberate: it measures the
architecture that a straightforward ROS~2 adoption produces.

\section{Experimental Configurations}
\label{sec:configurations}

A small number of configurations is used, chosen so that they represent
meaningful architectural differences rather than arbitrary divisions of
the pipeline:

\begin{itemize}
  \item \textbf{Configuration A --- Monolithic:} one Python process.
  \item \textbf{Configuration B --- Low partitioning:} two ROS~2 nodes.
  \item \textbf{Configuration C --- Higher partitioning:} four ROS~2
        nodes.
\end{itemize}

\todo{SPECIFY: the exact node boundaries for Configurations B and C, and
the justification that they represent meaningful architectural
differences.}

The primary experimental matrix crosses the three architectural
configurations with three input resolutions. Resolution is included as a
secondary dimension because both terms of the trade-off under study scale
with it: inference cost grows with input resolution \cite{szegedi2026},
and so does the communication cost of the image messages themselves,
since per-message serialization and copying costs grow with payload size
\cite{kronauer2021, macenski2022, luo2025}. The resolutions are selected
based on the capabilities of the chosen input source and model:

\begin{itemize}
  \item Resolution 1: \todo{SPECIFY (candidate: 320$\times$240)}
  \item Resolution 2: \todo{SPECIFY (candidate: 640$\times$480)}
  \item Resolution 3: \todo{SPECIFY (candidate: 1280$\times$720)}
\end{itemize}

This yields $3 \times 3 = 9$ experimental conditions:

% NOTE: requires \usepackage{amssymb} in the preamble for \checkmark.
\begin{table}[htbp]
  \centering
  \caption{Primary experimental matrix.}
  \label{tab:experimental-matrix}
  \begin{tabular}{clcccc}
    \hline
    \textbf{Config.} & \textbf{Architecture} & \textbf{Partitioning} &
    \textbf{Res.\ 1} & \textbf{Res.\ 2} & \textbf{Res.\ 3} \\
    \hline
    A & Single process & 1 process & \checkmark & \checkmark & \checkmark \\
    B & ROS 2          & 2 nodes   & \checkmark & \checkmark & \checkmark \\
    C & ROS 2          & 4 nodes   & \checkmark & \checkmark & \checkmark \\
    \hline
  \end{tabular}
\end{table}

Each condition is repeated multiple times, or executed for a sufficiently
long fixed duration, to obtain stable measurements.
\todo{SPECIFY: number of repetitions and run duration per condition.}

\section{Benchmarking Procedure}
\label{sec:procedure}

Each experimental condition follows the same procedure:

\begin{enumerate}
  \item \textbf{Initialize the system.} Start the application and all
        monitoring tools.
  \item \textbf{Warm-up.} Discard an initial number of frames or a fixed
        warm-up period to avoid measuring startup behavior: first-request
        costs such as model loading and memory allocation do not reflect
        steady-state operation.
        \todo{SPECIFY: warm-up duration or frame count.}
  \item \textbf{Execute the workload.} Run the pipeline for a fixed duration
        or a fixed number of frames.
        \todo{SPECIFY.}
  \item \textbf{Record measurements.} Collect end-to-end latency, FPS, CPU,
        RAM, and GPU utilization, and power consumption if measured.
  \item \textbf{Repeat.} Repeat the experiment sufficiently to characterize
        variability.
  \item \textbf{Reset.} Return the system to the same baseline state before
        the next experimental condition.
        \todo{SPECIFY: the reset procedure (e.g., stopping all nodes,
        cooling interval).}
\end{enumerate}

Throughout the measurements, the system is kept in a controlled state:
the fixed power mode is enforced, and background load is minimized. This
follows the practice of the ROS~2 measurement literature, where adaptive
operating-system behavior---CPU frequency scaling and idle states---is
explicitly neutralized because it adds noise to latency measurements
\cite{kronauer2021}.

Pilot experiments precede the final experiments in order to identify
measurement problems, unexpected bottlenecks, incorrect assumptions, and
configuration errors. The purpose of the pilot runs is to fix the
methodology, not to generate final results.

\section{Measurement of Performance Metrics}
\label{sec:metric-measurement}

\subsection{End-to-End Latency}

Latency is measured per frame as the elapsed time between the frame
entering the first pipeline stage and its fully postprocessed output being
available. Timestamps are taken at equivalent points in all
configurations---equivalent to the boundaries of the pipeline as a whole,
not of individual internal stages---so that measurements are comparable
across architectures. This mirrors the end-to-end perspective of
application-level ROS~2 latency evaluations \cite{betz2023a} and the
boundary-timestamping methodology used for layer-resolved communication
analysis \cite{kamiya2025}.
\todo{SPECIFY: timestamping mechanism (e.g., monotonic clocks) and, for
the ROS 2 configurations, how the timestamps of the first and final nodes
are related on the same host.} Reported statistics: mean, median, standard
deviation, minimum, maximum, P95, and P99 where sufficient samples are
available, following the percentile-reporting practice of the ROS~2
benchmarking literature \cite{kronauer2021, macenski2022}.

\subsection{Throughput}

Throughput is measured in frames per second over the steady-state
measurement interval, after warm-up removal:

\begin{equation}
  \mathrm{FPS} = \frac{\text{processed frames}}{\text{elapsed time}}
  \label{eq:fps}
\end{equation}

\subsection{CPU Utilization}

Average and peak CPU utilization are measured over the measurement
interval using system monitoring tools.
\todo{SPECIFY: tool and sampling interval (e.g., \texttt{tegrastats} or
equivalent).}

\subsection{Memory Utilization}

System RAM utilization is measured over the measurement interval; GPU
memory usage is recorded if reliable measurement is available.
\todo{SPECIFY: tool.}

\subsection{GPU Utilization}

GPU utilization is recorded to determine whether architectural overhead
changes the utilization of the primary compute resource.
\todo{SPECIFY: tool.}

\subsection{Power Consumption (Optional)}

If reliable measurement is available: average power, peak power, energy
per processed frame, and performance per watt, following the
power-metric practice of Orin Nano characterization studies
\cite{zequn2026}.
\todo{SPECIFY: measurement method, or state that power measurement was
excluded. Power measurement is an optional extension and must not
threaten completion of the thesis.}

\section{Data Analysis}
\label{sec:analysis}

For each experimental condition, the analysis proceeds as follows:

\begin{enumerate}
  \item Remove startup and warm-up measurements.
  \item Compute descriptive statistics for latency, throughput, and
        resource utilization.
  \item Compare latency distributions between configurations.
  \item Compare throughput and resource utilization between
        configurations.
  \item Compute percentage differences between configurations, including
        the relative overhead of each ROS~2 configuration with respect to
        the monolithic baseline:
        \begin{equation}
          \mathrm{overhead}\ (\%) = \frac{L_{\mathrm{ROS}} -
          L_{\mathrm{mono}}}{L_{\mathrm{mono}}} \times 100
          \label{eq:overhead}
        \end{equation}
        where $L$ denotes the metric under comparison (e.g., mean
        end-to-end latency).
  \item Investigate whether increasing node partitioning produces a
        consistent trend across resolutions.
\end{enumerate}

Where appropriate, statistical tests are selected after examining the
collected data and its distribution.
\todo{SPECIFY: tests to be applied once the data distribution is known.}

The analysis evaluates the following hypotheses, stated here as
predictions rather than expectations that must be confirmed; results
contradicting one or more hypotheses remain valid and will be analyzed as
such:

% NOTE: The thesis plan lists the hypotheses in the proposal but does not
% assign them a dedicated section in the chapter structure; they are
% placed here because the analysis evaluates them. Move if a better
% location is chosen.
\begin{description}
  \item[H1 (Latency):] The ROS~2 implementations will exhibit greater
        end-to-end latency than the monolithic implementation due to
        communication and processing overhead introduced between nodes.
  \item[H2 (Throughput):] The monolithic implementation will achieve equal
        or greater throughput than the ROS~2 implementations under
        equivalent computational conditions.
  \item[H3 (Resource utilization):] The ROS~2 implementations will
        require additional CPU and memory resources compared with the
        monolithic implementation.
  \item[H4 (Partitioning):] Increasing the number of ROS~2 communication
        boundaries will increase communication overhead and therefore
        increase end-to-end latency.
\end{description}

These predictions are grounded in the related literature: communication
overhead between processes has been measured repeatedly
\cite{kronauer2021, macenski2022, kamiya2025}, and mitigations that
reduce the number of process boundaries recover substantial resources
\cite{macenski2023}. The hypotheses nevertheless require experimental
verification for this workload, platform, and client library, since the
magnitude of the overhead relative to the dominant inference computation
is exactly what prior studies do not isolate.

\section{Reproducibility}
\label{sec:reproducibility}

The following information is documented so that the experiments can be
reproduced, in line with the community's call for standardized and
reproducible robotics benchmarking \cite{mayoral-vilches2024}:

\begin{itemize}
  \item Hardware model and configuration
  \item JetPack version
  \item Linux distribution
  \item ROS~2 distribution
  \item DDS implementation
  \item Python version
  \item Package versions
  \item Model version and weights
  \item Model configuration and precision
  \item Input resolution per experiment
  \item ROS~2 QoS configuration
  \item Power mode
  \item CPU/GPU configuration
  \item Ambient and environmental conditions, where relevant
  \item Benchmark duration
  \item Number of repetitions
  \item Warm-up procedure
\end{itemize}

Recording the full software and hardware configuration is not merely
administrative: latency evaluations of ROS~2 systems have shown that
results depend on the software and hardware configuration in non-obvious
ways \cite{betz2023a, kronauer2021}. The source code and benchmark
configuration are version-controlled. Any deviation from the planned
methodology observed during the final experiments is recorded and
reported.
